\section{Introduction}
\label{sec:intro}

Recently, a shift in the scaling law paradigm from pre-training to post-training has opened new possibilities for achieving another leap in AI model performance, as exemplified by the unprecedented AGI score of GPT-o3~\cite{openai:o3} and DeepSeek’s ``Aha moment''~\cite{Guo2025:deepseek}. Breakthroughs in LLMs have also extended to score-based generative models~\cite{sohldickstein2015deep, song2020sde, ho2020ddpm, albergo2023stochasticinterpolantsunifyingframework, ma2024sitexploringflowdiffusionbased}, resulting in significant improvements in user preference alignment~\cite{uehara2025inference}. Similar to the autoregressive generation process in LLMs, the denoising process in score-based generative models can be interpreted as a Sequential Monte Carlo (SMC)~\cite{doucet2001sequential, del2006sequential, chopin2020introduction} process with a single particle at each step. This perspective allows inference-time alignment to be applied analogously to LLMs by populating multiple particles at each step and selecting those that score highly under a given reward function~\cite{uehara2025inference, kim2025DAS, kim2025inference, wu2023tds, dou2024fps, singh2025code, li2024svdd}. A key distinction is that score-based generative models enable direct estimation of the final output from any noisy intermediate point via Tweedie’s formula~\cite{efron2011tweedie}, facilitating accurate approximation of the optimal value function~\cite{uehara2024understanding, uehara2025inference, uehara2024finetuningcontinuoustimediffusionmodels} through expected reward estimation. 

However, previous SMC-based approaches~\cite{kim2025DAS, dou2024fps, wu2023tds, trippe2023smcdiff, cardoso2024mcgdiff}, where each SMC step is coupled with the denoising process of score-based generative models, are limited in their ability to effectively explore high-reward regions, as the influence of the reward signal diminishes over time due to vanishing diffusion coefficient. Thus, rather than relying on particle exploration during later stages, it is more critical to identify effective initial latents that are well-aligned with the reward model from the outset.
In this work, we address this problem and propose an MCMC-based initial particle population method that generates strong starting points for the subsequent SMC process. 
This direction is particularly timely given recent advances in distillation techniques for score-based generative models~\cite{song2023consistency, kim2023consistency, lu2025scm, sauer2024ladd, liu2023flow}, now widely adopted in state-of-the-art models~\cite{BlackForestLabs:2024Flux, chen2025sanasprint}. These methods yield straighter generative trajectories and clearer Tweedie estimates~\cite{efron2011tweedie} from early steps, enabling more effective exploration from the reward-informed initial distribution.

A straightforward baseline for generating initial particles is the Top-$K$-of-$N$ strategy: drawing multiple samples from the standard Gaussian prior and selecting those with the highest reward scores. Though effective, this naive approach offers limited improvement in subsequent SMC due to its reliance on brute-force sampling. Motivated by these limitations, we explore Markov Chain Monte Carlo (MCMC)~\cite{metropolis_hastings, bj/1178291835, roberts2002langevin, duane1987hybrid, girolami2011riemann, Brooks_2011} methods based on Langevin dynamics, which are particularly well-suited to our setting since we sample from the initial posterior distribution, whose form is known. Nevertheless, applying MCMC in our problem presents unique challenges: the exploration space is extremely high-dimensional ($e.g.$, 65,536 for FLUX~\cite{BlackForestLabs:2024Flux}), posing significant challenges for conventional MCMC methods. In particular, the Metropolis–Hastings (MH) accept-reject mechanism, when used with standard Langevin-based samplers, becomes ineffective in such high-dimensional regimes, as the acceptance probability rapidly diminishes and most proposals are rejected.

Our key idea for enabling effective particle population from the initial reward-informed distribution is to leverage the Preconditioned Crank–Nicolson (pCN) algorithm~\cite{Cotter_2013pcnl, Beskos_2017geometricMCMC, doi:10.1142/S0219493708002378pcnlDiffusionBridge}, which is designed for function spaces or infinite-dimensional Hilbert spaces. When combined with the Langevin algorithm (yielding pCNL), its semi-implicit Euler formulation allows for efficient exploration in a high-dimensional space. Furthermore, when augmented with the MH correction, the acceptance rate is significantly improved compared to vanilla MALA. We therefore propose performing pCNL over the initial posterior distribution and selecting samples at uniform intervals along the resulting Markov chain. These samples are then used as initial particles for the subsequent SMC process across the denoising steps. We refer to the entire pipeline—\textbf{P}CNL-based initial particle sampling followed by \textbf{S}MC-based \textbf{I}nference-time reward alignment—as \textbf{PSI ($\Psi$)}-Sampler. To the best of our knowledge, this is the first work to apply the pCN algorithm in the context of generative modeling.

In our experiments, we evaluate three reward alignment tasks: layout-to-image generation (placing objects in designated bounding boxes within the image), quantity-aware generation (aligning the number of objects in the image with the specified count), and aesthetic-preference generation (enhancing visual appeal). We compare our \methodname{} against the base SMC method~\cite{wu2023tds} with random initial particle sampling, SMC combined with initial particle sampling via Top-$K$-of-$N$, ULA, and MALA, as well as single-particle methods~\cite{chung2023dps, yu2023freedom}. Across all tasks, \methodname{} consistently achieves the best performance in terms of the given reward and generalizes well to the held-out reward, matching or surpassing existing baselines. Its improvement over the base SMC method highlights the importance of posterior-based initialization, while its outperformance over ULA and MALA further confirms the limitations of these methods in extremely high-dimensional spaces.
